\documentclass[10pt]{article}
\usepackage[a4paper,margin=20mm]{geometry}
\usepackage{amsmath,amssymb,amsthm,mathtools,bm}
\usepackage{booktabs,array,microtype,xurl,hyperref}
\hypersetup{colorlinks=true,linkcolor=blue,citecolor=blue,urlcolor=blue}

\newtheorem{theorem}{Theorem}[section]
\newtheorem{proposition}[theorem]{Proposition}
\newtheorem{lemma}[theorem]{Lemma}
\newtheorem{corollary}[theorem]{Corollary}
\newtheorem{definition}[theorem]{Definition}
\newtheorem{remark}[theorem]{Remark}
\newcommand{\dd}{\,\mathrm d}
\newcommand{\RR}{\mathbb R}
\newcommand{\SSS}{\mathbb S}
\newcommand{\cE}{\mathcal E}
\newcommand{\cG}{\mathcal G}
\newcommand{\cV}{\mathcal V}
\newcommand{\Id}{\mathrm{Id}}

\title{AP-E17 Recovery-Compatible Boundary Transfer,\\
Strong-Quasiconvex Dirac Rigidity, and the\\
Nonvolumetric WZ-Capacity Alternative}
\author{Route-F derivation ledger}
\date{24 July 2026}

\begin{document}
\maketitle

\begin{abstract}
This note executes the three AP-E16 continuation gates and separates the
closed statements from the remaining endpoint assumption.  On the
volumetric part of the relaxation defect, a diagonal minimizing sequence
already has vanishing normalized \emph{full} local comparison deficit.  The
missing operation is not quasiminimality but replacement of its oscillatory
boundary trace by the affine barycentric tangent.

We formulate an explicit graph-tight annulus condition and prove a
recovery-compatible, sphere-valued boundary modification under precisely
that condition.  The comparison preserves the original trace, relative
homotopy, and global degree, while its normalized energy converges to the
affine tangent energy.  Combining it with the AP-E12 exact
strong-quasiconvexity identity gives an exact variance formula for the
homogeneous tangent graph Young measure.  Normalized quasiminimality then
forces both its ordinary variance and its concentration mass to vanish; the
measure is Dirac.  Consequently the volumetric defect is zero wherever the
graph-tight annulus condition holds.

The condition is not derived from the present endpoint graph bounds.  A
rank-two microball scaling shows why ordinary $W^{1,2}$ decomposition is
insufficient: strong $L^2$ convergence and bounded first/second-minor energy
can coexist with a divergent weighted gluing term.  Thus the result closes
the Young-measure rigidity implication, but not the endpoint annular
replacement theorem.

For every nonvolumetric tangent we prove an independent alternative.  The
local Wess--Zumino number $\mathsf b(r)$ obeys the sharp capacity inequality
\[
 \cE_3(B_r)\geq\frac{3K\pi^3}{2r^3}|\mathsf b(r)|^2 .
\]
Hence a defect of scale $r^d$ has
$|\mathsf b(r)|=O(r^{(d+3)/2})$ and cannot carry a nonzero quantized local
degree.  After normalization, either a positive fraction of the tangent is
forced by this vanishing but capacity-active WZ flux, or the tangent is
WZ-neutral and any remaining obstruction is analytic rather than
topological.  Hessian, determinant, and portal promotion remain closed.
\end{abstract}

\tableofcontents

\section{Frozen graph energy and the volumetric tangent}

The AP-E11 derivative density and complete-minor graph map are
\begin{align}
 W(F)&=\frac12|F|^2+\frac R2|M_2(F)|^2+\frac K2|M_3(F)|^2,
 & (R,K)&=(1,0.35),                                      \label{eq:W}\\
 \Xi(F)&=\bigl(F,\sqrt R\,M_2(F),\sqrt K\,M_3(F)\bigr),
 & W(F)&=\frac12|\Xi(F)|^2.                              \label{eq:Xi}
\end{align}
Let $u_j$ be the smooth degree-$B$ recovery sequence of AP-E15,
$\cE(u_j)\to I_B$, with derivative-energy defect
\begin{equation}
 W(Du_j)\dd x\stackrel{*}{\rightharpoonup}
 W(Du)\dd x+\mu .                                       \label{eq:defect}
\end{equation}

Define the volumetric set
\begin{equation}
 \cV=\left\{x_0:
 0<\vartheta(x_0)=
 \lim_{r\downarrow0}\frac{\mu(B_r(x_0))}{|B_r|}<\infty
 \right\}.                                               \label{eq:V}
\end{equation}
On $\cV$, $\mu$ is absolutely continuous with respect to Lebesgue measure.
We may therefore choose $x_0$ to be simultaneously a density point of
\eqref{eq:defect}, a Sobolev differentiability point of $u$, and a
localization point of all complete-minor Young-measure moments.  Put
\begin{equation}
 q=u(x_0),\qquad F_0=Du(x_0),\qquad F_0^Tq=0.             \label{eq:qF}
\end{equation}

Choose $r_\ell\downarrow0$ and then the diagonal $j_\ell$ so that
\begin{align}
 \frac{\cE(u_{j_\ell})-I_B}{{\mu(B_{r_\ell}(x_0))}}
 &\longrightarrow0,                                    \label{eq:global-diag}\\
 \frac1{r_\ell^5}\int_{B_{r_\ell}(x_0)}
 |u_{j_\ell}-u|^2\dd x&\longrightarrow0.                \label{eq:L2-diag}
\end{align}
The second choice is possible after the radii are fixed because
$u_j\to u$ strongly in $L^2$.  With $B=B_1(0)$, define the chordal tangent
\begin{equation}
 v_\ell(y)=\frac{u_{j_\ell}(x_0+r_\ell y)-q}{r_\ell},
 \qquad v_0(y)=F_0y.                                    \label{eq:v}
\end{equation}
Sobolev differentiability and \eqref{eq:L2-diag} give
\begin{equation}
 v_\ell\longrightarrow v_0\quad\hbox{strongly in }L^2(B).
 \label{eq:v-strong}
\end{equation}
Notice the exact scaling identity
$D_yv_\ell=Du_{j_\ell}(x_0+r_\ell y)$.

\section{What full local quasiminimality already gives}

For $Q\Subset\Omega$, AP-E16 defines the relative-homotopy local deficit
\begin{equation}
 \epsilon_j(Q)=\cE(u_j;Q)-
 \inf\{\cE(w;Q):w=u_j\text{ near }\partial Q,\
 w\simeq u_j\text{ rel. }\partial Q\}.                  \label{eq:eps}
\end{equation}
Replacing the map only in $Q$ preserves its exterior boundary value and
global degree, hence
\begin{equation}
 0\leq\epsilon_j(Q)\leq\cE(u_j)-I_B.                    \label{eq:eps-global}
\end{equation}

\begin{proposition}[Normalized full local quasiminimality]
For every fixed $0<\rho<1$, the diagonal
\eqref{eq:global-diag} satisfies
\begin{equation}
 \frac{\epsilon_{j_\ell}(B_{\rho r_\ell}(x_0))}
      {\mu(B_{r_\ell}(x_0))}\longrightarrow0.            \label{eq:normalized-eps}
\end{equation}
\end{proposition}

\begin{proof}
Apply \eqref{eq:eps-global} and divide by
$\mu(B_{r_\ell}(x_0))$.  This uses every target-valued, trace-preserving,
relative-homotopy comparison, not merely the compact fibre phase.
\end{proof}

Thus AP-E16's desired normalized quasiminimality was already contained in
the global diagonal.  The actual missing step is to build a comparison whose
trace is $u_{j_\ell}$ but whose interior is the affine tangent.

\section{Graph-tight annuli and boundary modification}

Let $A_\delta=B\setminus\overline{B_{1-\delta}}$.  Choose
$\eta_\ell\in C^\infty(B;[0,1])$ with
\begin{equation}
 \eta_\ell=0\text{ on }B_{1-2\delta_\ell},\qquad
 \eta_\ell=1\text{ near }\partial B,\qquad
 |D\eta_\ell|\leq C/\delta_\ell.                         \label{eq:eta}
\end{equation}

\begin{definition}[Graph-tight annulus condition, GTA]
The tangent sequence has a graph-tight annulus if there are
$\delta_\ell\downarrow0$ and $\eta_\ell$ as above such that
\begin{align}
 q\cdot u_{j_\ell}(x_0+r_\ell y)&\geq c_0>0
 \quad\text{on }A_{2\delta_\ell},                        \label{eq:hemisphere}\\
 \int_{A_{2\delta_\ell}}\!\!
 \bigl[W(Dv_\ell)+W(F_0)\bigr]\dd y&\longrightarrow0,    \label{eq:annular-energy}\\
 \int_{A_{2\delta_\ell}}\!\!
 \frac{|v_\ell-v_0|^2}{\delta_\ell^2}
 \bigl(1+R|Dv_\ell|^2+K|M_2(Dv_\ell)|^2\bigr)\dd y
 &\longrightarrow0.                                    \label{eq:annular-weight}
\end{align}
\end{definition}

Condition \eqref{eq:hemisphere} is the target-valued trace gate;
\eqref{eq:annular-energy} prevents energy from being trapped on the chosen
boundary layer; and \eqref{eq:annular-weight} is exactly the rank-one cutoff
remainder at all three minor orders.  No unspecified $o(1)$ is hidden in the
definition.

Set
\begin{equation}
 \widehat v_\ell
 =\eta_\ell v_\ell+(1-\eta_\ell)v_0,\qquad
 \widehat u_\ell(x_0+r_\ell y)
 =\frac{q+r_\ell\widehat v_\ell(y)}
 {|q+r_\ell\widehat v_\ell(y)|}.                        \label{eq:uhat}
\end{equation}
Outside $B_{r_\ell}(x_0)$ leave $\widehat u_\ell=u_{j_\ell}$.

\begin{lemma}[Rank-one graph estimate]
On the transition annulus,
\begin{equation}
 D\widehat v_\ell
 =\eta_\ell Dv_\ell+(1-\eta_\ell)F_0
 +(v_\ell-v_0)\otimes D\eta_\ell.                       \label{eq:Dvhat}
\end{equation}
The last term has rank one.  Consequently
\begin{align}
 W(D\widehat v_\ell)
 \leq C\bigg[&
 1+W(Dv_\ell)+W(F_0)\notag\\
 &+\frac{|v_\ell-v_0|^2}{\delta_\ell^2}
 \bigl(1+R|Dv_\ell|^2+K|M_2(Dv_\ell)|^2\bigr)\bigg].
 \label{eq:graph-cutoff}
\end{align}
\end{lemma}

\begin{proof}
The first-order estimate follows directly from \eqref{eq:Dvhat}.  If
$A=\eta_\ell Dv_\ell+(1-\eta_\ell)F_0$ and
$H=(v_\ell-v_0)\otimes D\eta_\ell$, then $\operatorname{rank}H\leq1$.
Multilinearity of determinants gives
\[
 M_2(A+H)=M_2(A)+L_2(A,H),\qquad
 M_3(A+H)=M_3(A)+L_3(A,A,H),
\]
with
$|L_2|^2\leq C|A|^2|H|^2$ and
$|L_3|^2\leq C|M_2(A)|^2|H|^2$.
Convex interpolation bounds the minors of $A$ by the corresponding
quantities for $Dv_\ell$ and $F_0$.  Using
$|H|\leq C|v_\ell-v_0|/\delta_\ell$ proves
\eqref{eq:graph-cutoff}.
\end{proof}

\begin{theorem}[Recovery-compatible boundary modification]
If GTA holds, then \eqref{eq:uhat} is well defined for all sufficiently
large $\ell$, is sphere valued, equals $u_{j_\ell}$ near
$\partial B_{r_\ell}(x_0)$, and is homotopic to it relative to that boundary.
Moreover
\begin{equation}
 \limsup_{\ell\to\infty}\frac1{r_\ell^3}
 \int_{B_{r_\ell}(x_0)}
 W(D\widehat u_\ell)\dd x
 \leq |B|W(F_0).                                       \label{eq:competitor-energy}
\end{equation}
The normalized difference of the potential terms also tends to zero.
\end{theorem}

\begin{proof}
Near the boundary, $\widehat v_\ell=v_\ell$, so
$q+r_\ell\widehat v_\ell=u_{j_\ell}$ and
$\widehat u_\ell=u_{j_\ell}$ exactly.  In the transition layer,
\eqref{eq:hemisphere} and $q+r_\ell v_0\to q$ uniformly keep the straight
interpolation away from the origin.  Radial projection therefore supplies a
relative-boundary homotopy and preserves global degree.

For $\mathcal P(z)=z/|z|$,
\begin{equation}
 D\mathcal P(z)=\frac{\Id-\mathcal P(z)\otimes\mathcal P(z)}{|z|}.
 \label{eq:DP}
\end{equation}
After the change of variables $x=x_0+r_\ell y$,
\[
 D_x\widehat u_\ell
 =D\mathcal P(q+r_\ell\widehat v_\ell)D_y\widehat v_\ell.
\]
On $B_{1-2\delta_\ell}$, $\widehat v_\ell=v_0$, and
$D_x\widehat u_\ell\to F_0$ uniformly.  On the annulus,
$D\mathcal P$ and its exterior powers are uniformly bounded.  Therefore
\eqref{eq:graph-cutoff} together with
\eqref{eq:annular-energy}--\eqref{eq:annular-weight} makes the annular
energy $o(1)$.  This proves \eqref{eq:competitor-energy}.

The potential $m^2(1-u^0)$ is Lipschitz on $\SSS^3$.  Both the affine
interior and the GTA transition converge in measure to $q$, while the
annular volume tends to zero.  Its normalized comparison error is therefore
$o(1)$.
\end{proof}

\section{Transfer to the tangent and exact Dirac rigidity}

\begin{theorem}[Volumetric quasiminimality transfer]
At a point $x_0\in\cV$ satisfying GTA,
\begin{equation}
 \limsup_{\ell\to\infty}\frac1{r_\ell^3}
 \int_{B_{r_\ell}(x_0)}W(Du_{j_\ell})\dd x
 \leq |B|W(F_0).                                       \label{eq:transfer}
\end{equation}
Consequently $\vartheta(x_0)=0$.
\end{theorem}

\begin{proof}
Use $\widehat u_\ell$ in \eqref{eq:eps} and combine
\eqref{eq:normalized-eps} with \eqref{eq:competitor-energy}.  Since
$\mu(B_{r_\ell})\sim\vartheta(x_0)|B|r_\ell^3$, the normalized local deficit
is $o(r_\ell^3)$.  On the other hand, differentiation of
\eqref{eq:defect} gives
\[
 \lim_{\ell\to\infty}\frac1{r_\ell^3}
 \int_{B_{r_\ell}(x_0)}W(Du_{j_\ell})\dd x
 =|B|\{W(F_0)+\vartheta(x_0)\}.
\]
Comparison with \eqref{eq:transfer} forces $\vartheta(x_0)=0$.
\end{proof}

The same conclusion has a more informative Young-measure form.  Let
$(\nu,\lambda,\nu^\infty)$ be the homogeneous generalized tangent Young
measure generated by the graph vectors $\Xi(Dv_\ell)$.  Here $\nu$ is the
ordinary oscillation measure, while $\lambda$ and $\nu^\infty$ record
quadratic graph concentration.  Localization and the null-Lagrangian moment
identities give
\begin{equation}
 \langle F\rangle_\nu=F_0,\qquad
 \langle M_k(F)\rangle_\nu=M_k(F_0),\quad k=2,3.
 \label{eq:moments}
\end{equation}

\begin{theorem}[Exact graph-variance identity and Dirac rigidity]
For the homogeneous tangent,
\begin{align}
 \overline W-W(F_0)
 ={}&\frac12\int\bigl[
 |F-F_0|^2
 +R|M_2(F)-M_2(F_0)|^2\notag\\
 &\hspace{35mm}
 +K|M_3(F)-M_3(F_0)|^2\bigr]\dd\nu(F)
 +\frac12\lambda(B).                                   \label{eq:variance}
\end{align}
If normalized quasiminimality is transferred by GTA, then
$\overline W\leq W(F_0)$ and every term on the right vanishes.  Hence
\begin{equation}
 \boxed{\nu=\delta_{F_0},\qquad\lambda=0.}               \label{eq:Dirac}
\end{equation}
\end{theorem}

\begin{proof}
Expand the three squares in $W$.  The cross terms vanish by
\eqref{eq:moments}; the quadratic recession of
$\frac12|\Xi|^2$ equals $\frac12$ on the graph unit sphere and contributes
$\lambda(B)/2$.  This proves \eqref{eq:variance}.  The transfer theorem gives
the reverse inequality, whereas \eqref{eq:variance} is nonnegative.  Equality
forces $F=F_0$ $\nu$-a.s. and zero concentration mass.
\end{proof}

This is the requested strong-quasiconvex conclusion: once the correct
boundary comparison exists, there is no nontrivial homogeneous tangent
Young measure left to classify.

\section{Why GTA is not yet a consequence of the graph bound}

The classical decomposition lemma for an $L^p$-bounded sequence of gradients
can separate an $L^p$-equiintegrable part from concentration
\cite{FonsecaMullerPedregal1998}.  Generalized gradient Young-measure theory
also describes boundary concentration for one fixed $p$-growth
\cite{Kruzik2012}.  Neither statement modifies one map while simultaneously
preserving the algebraic identities
\[
 M_2(Dv)=Dv\wedge Dv,\qquad M_3(Dv)=Dv\wedge Dv\wedge Dv,
\]
their separate endpoint $L^2$ bounds, the $\SSS^3$ target, and relative
degree.

The missing weighted estimate cannot be recovered from strong $L^2$
convergence by H\"older alone.  Let $\psi\in C_c^\infty(B_1;\RR^2)$ have a
nonzero second minor, embed it in $\RR^4$, and set on a ball of radius $h$
\begin{equation}
 z_h(x)=a_h\psi(x/h),\qquad a_h=h^{1/4}.                 \label{eq:microball}
\end{equation}
Its derivative has rank at most two, and scaling gives
\begin{align}
 \int|z_h|^2&\asymp h^{7/2},&
 \int|Dz_h|^2&\asymp h^{3/2},\notag\\
 \int|M_2(Dz_h)|^2&\asymp1,&
 M_3(Dz_h)&=0.                                         \label{eq:micro-scales}
\end{align}
Yet if the cutoff thickness is $\delta=h$,
\begin{equation}
 \delta^{-2}\int|z_h|^2|M_2(Dz_h)|^2
 \asymp h^{-3/2}\longrightarrow\infty.                  \label{eq:micro-weight}
\end{equation}
Radial projection of $(1,z_h)$ to $\SSS^3$ preserves these leading powers.
This example does not prove that every annulus is bad; it proves the precise
logical point that the current strong $L^2$ and complete-minor energy bounds
do not imply \eqref{eq:annular-weight}.  A complete-minor
equiintegrability/trace theorem, or a model-specific reverse-H\"older
estimate, is still required.

\section{Nonvolumetric concentration and WZ capacity}

For a smooth $u:B_r\to\SSS^3$, write
\begin{equation}
 J_u=\det[u,\partial_1u,\partial_2u,\partial_3u],\qquad
 \mathsf b_u(r)=\frac1{2\pi^2}\int_{B_r}J_u\dd x.        \label{eq:b}
\end{equation}
With the AP-E14 orientation, $\mathsf b$ is also
$(16\pi^2)^{-1}\int a\wedge da$.  Pointwise
$|J_u|\leq|M_3(Du)|$.

\begin{theorem}[Sharp local WZ-capacity lower bound]
For every ball,
\begin{equation}
 \boxed{\;
 \cE_3(u;B_r)=\frac K2\int_{B_r}|M_3(Du)|^2\dd x
 \geq\frac{3K\pi^3}{2r^3}|\mathsf b_u(r)|^2 .\;}
 \label{eq:wz-capacity}
\end{equation}
In Hopf variables the first/curvature components also obey
\begin{equation}
 \left(\frac18\int_{B_r}|a|^2\right)
 \left(\frac R{32}\int_{B_r}|da|^2\right)
 \geq R\pi^4|\mathsf b_u(r)|^2.                         \label{eq:hopf-capacity}
\end{equation}
\end{theorem}

\begin{proof}
Cauchy--Schwarz and $|B_r|=4\pi r^3/3$ give
\[
 (2\pi^2|\mathsf b|)^2
 \leq |B_r|\int_{B_r}|M_3(Du)|^2.
\]
Multiplication by $K/2$ proves \eqref{eq:wz-capacity}.
For \eqref{eq:hopf-capacity}, use
$(16\pi^2\mathsf b)^2\leq
(\int|a|^2)(\int|da|^2)$.
\end{proof}

\begin{corollary}[No quantized point-degree concentration]
Suppose along a tangent sequence
\begin{equation}
 \cE(u_{j_\ell};B_{r_\ell}(x_0))\leq Cr_\ell^d
 \quad(d\geq0).
 \label{eq:dscale}
\end{equation}
Then
\begin{equation}
 |\mathsf b_{u_{j_\ell}}(r_\ell)|
 \leq\left(\frac{2C}{3K\pi^3}\right)^{1/2}
 r_\ell^{(d+3)/2}\longrightarrow0.                     \label{eq:bdecay}
\end{equation}
If the boundary is constant so that $\mathsf b$ is an integer relative
degree, it is therefore zero for all sufficiently large $\ell$.
\end{corollary}

Thus no nonvolumetric defect can hide a nonzero integer $B=1$ bubble at a
point.  The Hopf base has no separate monopole charge either: since $a$ is a
globally defined one-form on the ball,
\begin{equation}
 \int_{\partial B_r}da=\int_{B_r}d(da)=0.               \label{eq:no-monopole}
\end{equation}

For rational $0<\rho<1$, define the normalized capacity fraction
\begin{equation}
 \kappa_\ell(\rho)=
 \frac{3K\pi^3}{2}
 \frac{|\mathsf b_{u_{j_\ell}}(\rho r_\ell)|^2}
 {(\rho r_\ell)^3\mu(B_{r_\ell}(x_0))}.                 \label{eq:kappa}
\end{equation}
After one diagonal subsequence, every $\kappa_\ell(\rho)$ has a limit.
Equation \eqref{eq:wz-capacity} gives the following exhaustive alternative.

\begin{theorem}[Nonvolumetric capacity/topology alternative]
At every nonvolumetric tangent point, exactly one of the following occurs.
\begin{enumerate}
 \item \textbf{Capacity-active branch.}  For some rational $\rho$,
 $\lim\kappa_\ell(\rho)>0$.  A positive fraction of the normalized tangent
 energy is forced into the sextic/WZ component.  Its absolute WZ number still
 tends to zero by \eqref{eq:bdecay}, so it is a fractional capacity effect,
 not a quantized topological atom.
 \item \textbf{WZ-neutral branch.}  For every rational $\rho$,
 $\kappa_\ell(\rho)\to0$.  The tangent carries neither an integer local
 degree nor a positive normalized WZ-capacity fraction.  Any surviving
 defect is analytically concentrated in the graph variables and is not
 protected by Route-E topology.
\end{enumerate}
\label{thm:alternative}
\end{theorem}

The alternative is a classification, not yet an exclusion theorem.  The
capacity-active branch must be retained as an order-specific concentration
coordinate.  The WZ-neutral branch can be removed only after a GTA-type
trace theorem is proved for generalized graph tangents.

\section{Gate decision}

The three requested tasks now have the following exact status.
\begin{enumerate}
 \item \textbf{Boundary transfer: conditional theorem proved.}
 Normalized full local quasiminimality is unconditional on the AP-E16
 diagonal.  Its transfer to the affine tangent is proved under GTA.  GTA
 itself is not implied by the present endpoint graph bounds.
 \item \textbf{Dirac rigidity: proved once transfer holds.}
 Equation \eqref{eq:variance} forces both oscillation variance and
 concentration mass to vanish, so the homogeneous tangent is exactly
 $\delta_{F_0}$.
 \item \textbf{Nonvolumetric alternative: proved unconditionally.}
 Equations \eqref{eq:wz-capacity}--\eqref{eq:kappa} exclude quantized
 point-degree concentration and split every remaining tangent into the
 capacity-active or WZ-neutral branch.
\end{enumerate}

The new mainline blocker is therefore narrower than AP-E16: prove GTA for
the actual minimizing diagonal, most plausibly by a complete-minor
equiintegrability theorem or a model-specific reverse-H\"older estimate.
Until that endpoint trace theorem closes, full $\mu=0$, local recovery,
classicality, isolation, and the bosonic Hessian remain false.  Determinant
and degree-one portal promotion also remain false.

\paragraph{Reproducibility.}
Run \texttt{python3}
\path{route_f/code/verify_ap_e17_boundary_transfer_dirac_capacity.py}
from the repository root.  The card verifies radial projection, exact
strong-quasiconvex variance, the microball scaling obstruction, WZ-capacity
constants, the nonvolumetric exponent, and every fail-closed gate.  It
performs no lattice relaxation or grid scan.

\begin{thebibliography}{9}
\bibitem{FonsecaMullerPedregal1998}
I.~Fonseca, S.~M\"uller, and P.~Pedregal,
\emph{Analysis of concentration and oscillation effects generated by
gradients}, SIAM J. Math. Anal. \textbf{29} (1998), 736--756,
\url{https://doi.org/10.1137/S0036141096306534}.

\bibitem{Kruzik2012}
M.~K\v{r}u\v{z}\'ik,
\emph{Quasiconvexity at the boundary and concentration effects generated by
gradients}, ESAIM Control Optim. Calc. Var. \textbf{19} (2013), 679--700,
\url{https://arxiv.org/abs/1009.0795}.
\end{thebibliography}

\end{document}
